% CCSS 7.SP.C.7 — Probability Models
\section{Probability Models}

\begin{learningGoals}
\begin{itemize}
    \item Develop and compare uniform and non-uniform probability models
    \item Compare predicted probabilities to observed frequencies
\end{itemize}
\end{learningGoals}

\begin{conceptBox}{Probability Models}
\begin{itemize}[leftmargin=*]
    \item \textbf{Uniform model}: All outcomes equally likely
    \item \textbf{Non-uniform model}: Outcomes have different probabilities
\end{itemize}
\end{conceptBox}

\begin{workedExample}{Spinner Model}
A spinner has $4$ equal sections. What is the probability of landing on any section?

Uniform: $1/4$
\answerHighlight{$1/4$}
\end{workedExample}

\tipBox{Compare predictions to actual results to check your model.}

\begin{practiceBox}{Probability Models}
\resetProblems

\prob A bag has $2$ red, $3$ blue, $5$ green marbles. Is this a uniform model?
\answerExplain{No}{Probabilities are different for each color.}

\prob A die is rolled. Is this a uniform model?
\answerExplain{Yes}{Each number $1$–$6$ is equally likely.}

\prob A spinner has $3$ sections: $1$ large, $2$ small. Is this uniform?
\answerExplain{No}{Sections are not equal, so probabilities differ.}

\prob If a coin is weighted to land heads $70\%$ of the time, is this uniform?
\answerExplain{No}{Probabilities are not equal for heads and tails.}

\prob Why compare predicted probabilities to observed frequencies?
\answerExplain{To check model accuracy}{If results differ, the model may need adjustment.}

\end{practiceBox}
